% META: {"id":"1SPE-DERGLOBAL-CO-051","chapitre":"1SPE-DERIVATION-GLOBAL","type_objet":"corrige","exercice_id":"1SPE-DERGLOBAL-EX-051","status":"approved","origin":"generated","status_history":["generated","audited","accepted","approved"],"release_acceptance":"RELEASE_OWNER_BATCH_ACCEPTANCE","acceptance_closure_digest":"sha256:9b3ccf9a81c5520fbb7e03b7b2d3e7bf2057a02834b6d908bf3be5f49f3b553f","acceptance_receipt_digest":"sha256:4fad86f0282e0fa4e0419cca1ad6379ae7af5a280c04a4a90880f90a1c044242"}
% BEGIN-VERIFY
% from sympy import symbols
% x = symbols('x')
% f = 3*x**2 - 5
% assert f.subs(x, -x) == f
% g = x**3 + 2*x
% assert g.subs(x, -x) == -g
% h = x**2 + x
% h_neg = h.subs(x, -x)  # x^2 - x
% assert h_neg != h and h_neg != -h
% END-VERIFY
\begin{corrige}{1SPE-DERGLOBAL-EX-051}
Chaque fonction est définie sur $\mathbb{R}$, domaine symétrique par rapport à $0$.

\textbf{1.} Calculons $f(-x)$ :
\[
f(-x) = 3(-x)^2 - 5 = 3x^2 - 5 = f(x).
\]
Donc $f$ est \textbf{paire}. Sa courbe est symétrique par rapport à l'axe des ordonnées.

\textbf{2.} Calculons $g(-x)$ :
\[
g(-x) = (-x)^3 + 2(-x) = -x^3 - 2x = -(x^3 + 2x) = -g(x).
\]
Donc $g$ est \textbf{impaire}. Sa courbe est symétrique par rapport à l'origine.

\textbf{3.} Calculons $h(-x)$ :
\[
h(-x) = (-x)^2 + (-x) = x^2 - x.
\]
On a $h(-x) = x^2 - x
\neq x^2 + x = h(x)$ et $h(-x) = x^2 - x
\neq -(x^2 + x) = -x^2 - x = -h(x)$.

Donc $h$ n'est \textbf{ni paire ni impaire}.
\end{corrige}
